\chapter{Opis założeń projektu}
\label{cha:zalozenia}

Celem projektu jest stworzenie aplikacji desktopowej wspierającej nauczyciela w tworzeniu
i przeprowadzaniu testów oraz ucznia w ich rozwiązywaniu. Tradycyjne, papierowe sprawdzanie wiedzy
jest czasochłonne - wymaga ręcznego przygotowania arkuszy, ich powielenia, a następnie żmudnego
sprawdzania i podliczania punktów. Dodatkowo gromadzenie wyników i analiza, które zagadnienia
sprawiają uczniom najwięcej trudności, są w takim modelu praktycznie niemożliwe do przeprowadzenia
na bieżąco.

Aplikacja \emph{Testownik} rozwiązuje te problemy, przenosząc cały proces do jednego, spójnego
narzędzia. Nauczyciel układa test w edytorze, przypisuje go wybranym grupom uczniów i otrzymuje
automatycznie policzone wyniki wraz ze statystykami. Uczeń loguje się do systemu, rozwiązuje
przypisane mu testy w ograniczonym czasie i natychmiast poznaje swój wynik. Realizacja projektu
wymagała znajomości języka C\# i programowania obiektowego, technologii WPF do budowy interfejsu
graficznego, biblioteki Entity Framework Core do obsługi bazy danych, projektowania relacyjnych baz
danych PostgreSQL oraz systemu kontroli wersji Git.

\section{Wymagania funkcjonalne projektu}
\label{sec:wymaganiaFunkcjonalne}

\begin{itemize}
  \item \textbf{Zarządzanie kontami i rolami użytkowników}
  \begin{itemize}
    \item Użytkownik może samodzielnie zarejestrować nowe konto (imię, nazwisko, login, hasło),
          wybierając przy rejestracji rolę - ucznia lub nauczyciela - oraz zalogować się do
          systemu za pomocą loginu i hasła. Login jest unikalny w skali całego systemu.
    \item System rozróżnia dwie role - nauczyciela i ucznia - i po zalogowaniu udostępnia każdej
          z nich odrębny zestaw funkcji oraz inny układ nawigacji. Każdy nauczyciel zarządza
          wyłącznie własnymi testami.
  \end{itemize}

  \item \textbf{Zarządzanie grupami uczniów}
  \begin{itemize}
    \item Nauczyciel może tworzyć nowe grupy (np. klasy) opatrzone nazwą i opisem.
    \item Nauczyciel może przypisywać uczniów do grup oraz usuwać ich z grup; jeden uczeń może
          należeć do wielu grup jednocześnie.
    \item System wyświetla skład każdej grupy oraz przypisane do niej testy.
  \end{itemize}

  \item \textbf{Tworzenie i edycja testów}
  \begin{itemize}
    \item Nauczyciel tworzy testy złożone z pytań trzech typów: jednokrotnego wyboru, wielokrotnego
          wyboru oraz pytań otwartych.
    \item Dla każdego testu określa nazwę, przedmiot (wpisywany dowolnie lub wybierany z listy),
          limit czasu, próg zaliczenia (w procentach), dozwoloną liczbę podejść oraz grupy, którym
          test jest udostępniany.
    \item Dla każdego pytania określa treść, liczbę punktów oraz - w pytaniach zamkniętych -
          warianty odpowiedzi i poprawne odpowiedzi, a w pytaniach otwartych odpowiedź wzorcową.
    \item Test ma stan widoczności - aktywny lub nieaktywny. Tylko aktywny test jest widoczny dla
          uczniów; nowy test powstaje jako nieaktywny (szkic), a nauczyciel aktywuje go, gdy jest
          gotowy do przeprowadzenia.
    \item Test można edytować, duplikować oraz usuwać. Jeżeli test ma już rozwiązania uczniów,
          edycja pytań jest blokowana (edytowalne pozostają ustawienia i grupy), aby nie naruszyć
          powiązanych z pytaniami odpowiedzi.
  \end{itemize}

  \item \textbf{Rozwiązywanie testów przez ucznia}
  \begin{itemize}
    \item Uczeń widzi listę testów przypisanych do jego grup, z podziałem na rozwiązane
          i nierozwiązane.
    \item Podczas rozwiązywania działa licznik czasu; po jego upływie test jest automatycznie
          kończony i oddawany.
    \item Uczeń może swobodnie przechodzić między pytaniami, a przy oddawaniu testu z pominiętymi
          pytaniami otrzymuje stosowne ostrzeżenie.
    \item Liczba podejść do testu jest ograniczona wartością ustawioną przez nauczyciela.
  \end{itemize}

  \item \textbf{Ocenianie odpowiedzi}
  \begin{itemize}
    \item Pytania zamknięte są oceniane automatycznie - w pytaniach wielokrotnego wyboru
          z uwzględnieniem punktów częściowych.
    \item Pytania otwarte trafiają do kolejki oceniania, gdzie nauczyciel przyznaje punkty ręcznie,
          po czym wynik ucznia jest automatycznie przeliczany.
    \item Nauczyciel może ręcznie edytować ocenę dowolnego podejścia - przyznać lub nadpisać punkty
          każdego pytania, także automatycznie ocenionego pytania zamkniętego; wynik jest wówczas
          ponownie przeliczany.
    \item Przy wielu podejściach jako ostateczny brany jest najlepszy wynik ucznia.
  \end{itemize}

  \item \textbf{Wykrywanie prób oszustwa}
  \begin{itemize}
    \item System wykrywa opuszczenie okna testu (np. przełączenie się na inną aplikację). Pierwsze
          opuszczenie skutkuje ostrzeżeniem, a kolejne - automatycznym zakończeniem testu
          z wynikiem 0\% i oznaczeniem próby oszustwa.
  \end{itemize}

  \item \textbf{Statystyki i eksport wyników}
  \begin{itemize}
    \item Nauczyciel ma dostęp do statystyk testu: zdawalności, średniego wyniku, najtrudniejszego
          pytania oraz poprawności w rozbiciu na poszczególne pytania.
    \item Wyniki testu można wyeksportować do pliku CSV.
  \end{itemize}
\end{itemize}

\section{Wymagania niefunkcjonalne projektu}
\label{sec:wymaganiaNiefunkcjonalne}

System spełnia następujące kryteria jakościowe i ograniczenia:

\begin{itemize}
  \item \textbf{Użyteczność} - interfejs użytkownika jest spójny i intuicyjny. Nawigacja odbywa
        się przez stały panel boczny, a poszczególne ekrany utrzymane są w jednolitej stylistyce
        (karty, znaczniki statusu, wykresy), co ułatwia poruszanie się po aplikacji.

  \item \textbf{Niezawodność i obsługa błędów} - operacje na bazie danych oraz operacje
        wejścia/wyjścia są zabezpieczone wspólnym mechanizmem obsługi wyjątków. W razie błędu
        użytkownik otrzymuje zrozumiały komunikat, a aplikacja nie ulega awarii. Dane wprowadzane
        w formularzach (np. ustawienia testu, kompletność pytań) są walidowane przed zapisem.

  \item \textbf{Wydajność} - kluczowe operacje, takie jak logowanie, wczytywanie testów czy zapis
        wyniku, wykonują się w czasie nieprzekraczającym kilku sekund, a uruchomienie aplikacji
        nie przekracza kilku sekund.

  \item \textbf{Bezpieczeństwo} - dostęp do funkcji systemu jest możliwy wyłącznie po zalogowaniu,
        a zakres uprawnień zależy od roli użytkownika. Hasła nie są przechowywane jawnie - zapisuje
        się ich skrót (SHA-256). Unikalność loginów jest wymuszana na poziomie bazy danych.

  \item \textbf{Środowisko} - aplikacja jest przeznaczona na system Windows z zainstalowanym
        środowiskiem uruchomieniowym .NET oraz dostępem do serwera bazy danych PostgreSQL.

  \item \textbf{Integralność danych} - spójność danych zapewniają klucze obce wraz z odpowiednio
        dobranymi regułami usuwania (kaskadowymi i ograniczającymi) oraz mechanizm migracji
        Entity Framework Core, wersjonujący schemat bazy danych.

  \item \textbf{Utrzymywalność} - kod źródłowy został napisany zgodnie z zasadami programowania
        obiektowego i podzielony na warstwy (model, dostęp do danych, logika/usługi, interfejs
        użytkownika), co ułatwia jego rozwój i modyfikację.
\end{itemize}
